%% Введение

Полуструктурированные данные в формате JSON прочно заняли место
основного способа обмена данными между веб-сервисами, форматом
логирования событий, телеметрии устройств и хранилищем данных
социальных сетей. Одновременно традиционные реляционные СУБД
оказались плохо приспособлены к такого рода нагрузкам: хранение
JSON в виде неделимого значения требует полного считывания
документа с диска при обращении к единственному полю, а
классическое колоночное хранение нацелено на заранее известную
фиксированную схему.

Современные промышленные СУБД решают эту задачу по-разному.
Документоориентированные системы (MongoDB как исторический
законодатель этого направления, Couchbase) оптимизируют работу с
произвольной структурой за счёт отказа от реляционного интерфейса
и колоночного хранения. Колоночные аналитические СУБД (ClickHouse,
Snowflake, DuckDB) вводят специализированные типы данных,
сочетающие заранее объявленные колонки и общую область хранения
редких полей. Каждый из подходов сопряжён с компромиссами:
документные системы не обеспечивают производительности
аналитических запросов, а специализированные колоночные форматы
ограничены в выборе операций над редкими полями.

Настоящая работа посвящена разработке альтернативного механизма
хранения полуструктурированных данных в рамках колоночной СУБД
Otterbrix~--- проекта с открытым исходным кодом, развивающегося как
встраиваемый аналитический движок с поддержкой SQL. Отличительная
особенность предлагаемого решения заключается в сочетании колоночного
хранения частых полей и размещения редких полей в отдельных
вспомогательных двухколоночных таблицах, связанных с основной через
логический идентификатор. Такая архитектура даёт возможность
переиспользовать всю механику колоночного движка (векторизация,
zone~maps, индексы, MVCC, WAL) для редких полей без необходимости
поддерживать отдельную подсистему хранения.

В рамках работы были решены следующие задачи. Во-первых, выполнен
обзор существующих подходов к хранению JSON в СУБД: BSON в
MongoDB как канонический пример документной модели, модель
Entity-Attribute-Value и её варианты, колоночное развёртывание
путей, новый тип \texttt{JSON} в ClickHouse. Во-вторых, предложена
модель гетерогенного хранения, в которой статус поля (закреплённое
в основной таблице или вынесенное в отдельную side-таблицу)
фиксируется при \texttt{CREATE TABLE}, а множество полей таблицы
динамически расширяется по мере появления новых ключей в
поступающих JSON-документах. В-третьих, модель реализована в
кодовой базе Otterbrix: введён класс \texttt{computed\_schema},
инкапсулирующий динамическую схему, добавлены четыре физических
оператора (\texttt{scan\_computing\_table} для гибридного
сканирования и три варианта DML \texttt{*\_computing} для
multi-collection-операций под одной транзакцией), pre-flight в
диспетчере для DDL side-таблиц, поддержка SQL-конструкций
\texttt{CREATE TABLE db.t ();},
\texttt{WITH (pinned\_columns=\ldots)} и
\texttt{INSERT \ldots\ VALUES (json('\ldots'))}. В-четвёртых,
выполнены оптимизации физических операторов: проекция колонок на
уровне оптимизатора, типизированные предикаты фильтрации с
пакетным копированием подошедших строк, специализированные пути
агрегации, исправление ошибки в реализации zone~maps. В-пятых,
проведено сравнительное тестирование на стандартном бенчмарке
JSONBench: на трёх конфигурациях самой Otterbrix (исходная
документная, новая колоночно-разреженная, новая
колоночно-разреженная с оптимизациями) и в сравнении с
встраиваемым колоночным движком DuckDB.

Результаты показали, что архитектурный переход от документной
модели к колоночно-разреженной даёт ускорение в 3--14~раз;
оптимизации физических операторов добавляют ещё 4--23~раза.
Финальная конфигурация Otterbrix на запросах JSONBench
сопоставима с DuckDB на аналитических запросах по плотным полям
(отставание единицы процентов) и обходит DuckDB в 1.3--3.6~раз на
запросах с полным сканированием и обращением к разреженным полям.
Объём новой кодовой базы измеряется тысячами строк~--- на порядок
меньше того, что потребовался бы для разработки специализированной
подсистемы хранения с нуля.

Настоящая работа организована следующим образом. В главе~1
изложена постановка задачи, описана исходная архитектура Otterbrix
и мотивация выбранного направления. Глава~2 содержит обзор
существующих подходов. В главе~3 описывается спроектированная
модель хранения. Глава~4 посвящена деталям реализации. В главе~5
рассматриваются оптимизации физических операторов. Глава~6
приводит результаты сравнительного тестирования с DuckDB и с
тремя конфигурациями самой Otterbrix.
